\section{Route evaluation and source boundary}

The strongest positive result is an exact localization: a common aggregate
phase is absent from the same-block principal covariance, while every
nonorthogonal sign effect lies on a symmetrized cross-block edge.  TPC-243 then
confines physical hard-window variation to two Gram-error terms.

The strongest obstruction is equally explicit.  The local principal term is
\[
|C_h|^2\ip{w_h}{b_h},
\]
so the outer sign of $C_h$ cannot drive its cancellation.  A successful signed
mechanism must instead use at least one of the following: cancellation inside
$|C_h|$, a nontrivial local covariance, cross-block leakage, or asymmetric
multipliers on the two lanes.

The first physical blocker is the absent literal V59 map
\[
(\beta,w)\longmapsto\{b_{h,a},w_{h,a}\}_{(a,h)=1}
\]
with the same coefficient placement and one synthesis operator.  The current
result is therefore conditional at that interface.  It supplies no arithmetic
$L^2$, fixed-atom credit, strict endpoint payment, full Gate B, or twin-prime
conclusion.

The next minimal invariant is the within-block covariance.  Given a canonical
unit vector in a block, its longitudinal moments and transverse energies
determine a sharp center-radius disk for $\ip{w_h}{b_h}$.  This avoids chasing
an outer sign that the present theorem proves invisible.
